\section{Related Work}
\label{sec:related_work}

\noindent\textbf{Post-Harvest Vision and Post-Hoc OOD Detection:} Deep convolutional architectures have advanced automated produce sorting, defect identification, and ripeness estimation \cite{rizzo2023fruit,zhang2024fruit,chen2024classification,gulzar2023fruit}. Lightweight convolutional backbones such as MobileNetV2 \cite{sandler2018mobilenetv2} are frequently explored for embedded sorting applications \cite{chen2024classification,gulzar2023fruit} due to compact parameter footprints (approximately $2.26\times 10^6$ total parameters in our six-class configuration) and low computational overhead. However, closed-set classifiers cannot detect anomalous or non-target produce \cite{guo2017calibration}. Post-hoc OOD detection algorithms compute an anomaly score $S(x)$ from internal activations without retraining \cite{yang2024generalized}. Baseline Maximum Softmax Probability (MSP) \cite{hendrycks2017baseline} evaluates $\max_k p_k(x)$. Liang et al.\ \cite{liang2018enhancing} introduced ODIN, combining temperature scaling with input gradient perturbations: $\tilde{x} = x - \varepsilon \operatorname{sign}(-\nabla_x \log p_{\hat{y}}(x; T))$. Bendale and Boult \cite{bendale2016towards} formulated OpenMax by fitting Weibull distributions to penultimate features. Lee et al.\ \cite{lee2018simple} proposed class-conditional Gaussian modeling via Mahalanobis distance $d_M^2(x) = \min_k (f(x) - \mu_k)^\top \mathbf{\Sigma}^{-1} (f(x) - \mu_k)$. Liu et al.\ \cite{liu2020energy} proposed an energy-based OOD scoring function $S_{\text{Energy}}(x) = T \log \sum_{k=1}^K \exp(z_k(x)/T)$ computed over raw logits. The Energy score is a post-hoc logit-based OOD scoring function motivated by an energy-based formulation; it is not treated here as a calibrated estimate of the true input marginal density. Negated Euclidean Mahalanobis distance $-\sqrt{d_M^2(x)}$ strictly preserves sample ranking and AUROC while scaling linearly with feature deviations.

\vspace{2pt}
\noindent\textbf{Agricultural Domain Shifts and Research Positioning:} Recent studies explore domain adaptation and anomaly rejection in agricultural contexts. Wang et al.\ \cite{wang2023crossdomain} evaluated attention networks with unsupervised domain adaptation for fruit classification across lighting conditions. Pehlivan and Kara \cite{pehlivan2026lightweight} proposed one-shot open-set metric learning for food recognition. Jansi and Kavitha \cite{jansi2026peos} developed PEOS-Net, combining prototype learning and energy scores for crop disease recognition, while Yan et al.\ \cite{yan2026pagmsbnet} introduced PAG-MSBNet for open-set crop disease diagnosis. Despite these advances, many evaluations assess models on curated test splits collected under uniform protocols, which can leave camera hardware and background differences entangled with botanical distribution shifts. In contrast, this study evaluates a frozen standard MobileNetV2 classifier across multiple external repositories, establishing a source-stratified benchmark that systematically contrasts external covariate shift, semantic species shift, and dataset-reported chemical treatment shift under prospective, locked operating thresholds.
